\documentclass{article}
\usepackage{amsmath}
\usepackage{amssymb}

\begin{document}

\section{Model 3: Environment-Specific Ability Slopes with Task Random Effects}

\subsection{Model Specification}

Model 3 extends the baseline IRT framework by allowing environment-specific modulation of ability effects while incorporating task-level random effects. The model is specified as:

\begin{equation}
\text{logit } P(y_{ijkt} = 1) = \beta_0 + \lambda_{ij}^\text{eff} \cdot \theta_{ij}^K + \psi_{ij}^\text{eff} \cdot \theta_{ij}^R + \gamma_s + \delta_\ell + \xi_v + \alpha_j + \eta_t
\end{equation}

where the effective ability slopes are decomposed as:
\begin{align}
\lambda_{ij}^\text{eff} &= \lambda_\text{base} + \theta_j \\
\psi_{ij}^\text{eff} &= \psi_\text{base} + \phi_j
\end{align}

\subsection{Term Definitions}

\begin{itemize}
    \item $y_{ijkt}$: Binary success indicator for model $i$ in environment $j$ on task $t$ at level $k$

    \item $\beta_0$: Global intercept (overall baseline log-odds of success)

    \item \textbf{Knowledge ability component:}
    \begin{itemize}
        \item $\lambda_\text{base}$: Global knowledge loading (average effect of knowledge across all environments)
        \item $\theta_j$: Environment-specific adjustment to knowledge loading (deviation from global mean)
        \item $\theta_{ij}^K$: Model $i$'s knowledge ability in environment $j$ (estimated from knowledge Q\&A, standardized)
        \item $\lambda_{ij}^\text{eff} = \lambda_\text{base} + \theta_j$: Effective knowledge loading for environment $j$
    \end{itemize}

    \item \textbf{Reasoning ability component:}
    \begin{itemize}
        \item $\psi_\text{base}$: Global reasoning loading (average effect of reasoning across all environments)
        \item $\phi_j$: Environment-specific adjustment to reasoning loading (deviation from global mean)
        \item $\theta_{ij}^R$: Model $i$'s reasoning ability in environment $j$ (estimated from reasoning Q\&A, standardized)
        \item $\psi_{ij}^\text{eff} = \psi_\text{base} + \phi_j$: Effective reasoning loading for environment $j$
    \end{itemize}

    \item $\gamma_s$: Scaffold effect (e.g., ReAct vs Direct prompting)

    \item $\delta_\ell$: Difficulty level effect (ordered constraint: $\delta_1 > \delta_2 > \cdots > \delta_L$)

    \item $\xi_v$: Verbosity level effect (low, medium, high)

    \item $\alpha_j$: Environment main effect (captures baseline difficulty/characteristics of environment $j$)

    \item $\eta_t$: Task-specific random effect (captures unique characteristics of task $t$)
\end{itemize}

\subsection{Prior Distributions}

The hierarchical structure is completed with the following priors:

\begin{align}
\beta_0 &\sim \mathcal{N}(\text{logit}(\bar{y}), 10) \\
\lambda_\text{base} &\sim \mathcal{N}(0, 1) \\
\psi_\text{base} &\sim \mathcal{N}(0, 1) \\
\theta_j &\sim \mathcal{N}(0, 0.5) \quad \text{subject to } \sum_j \theta_j = 0 \\
\phi_j &\sim \mathcal{N}(0, 0.5) \quad \text{subject to } \sum_j \phi_j = 0 \\
\gamma_s &\sim \mathcal{N}(0, 1) \\
\delta_\ell &= -\sum_{m=1}^\ell \exp(\delta_\ell^\text{raw}), \quad \delta_\ell^\text{raw} \sim \mathcal{N}(0, 1) \\
\xi_v &\sim \mathcal{N}(0, 1) \\
\alpha_j &\sim \mathcal{N}(0, \sigma_\alpha), \quad \sigma_\alpha \sim \text{HalfNormal}(1) \\
\eta_t &\sim \mathcal{N}(0, \sigma_\eta), \quad \sigma_\eta \sim \text{HalfNormal}(0.5)
\end{align}

\subsection{Key Features}

\subsubsection{Sum-to-Zero Constraints}
The environment-specific adjustments $\theta_j$ and $\phi_j$ are constrained to sum to zero:
\begin{equation}
\sum_{j=1}^{J} \theta_j = 0, \quad \sum_{j=1}^{J} \phi_j = 0
\end{equation}

This constraint ensures:
\begin{enumerate}
    \item \textbf{Identifiability}: Separates the global effect ($\lambda_\text{base}$, $\psi_\text{base}$) from environment-specific deviations
    \item \textbf{Interpretability}: $\lambda_\text{base}$ represents the average knowledge effect; positive $\theta_j$ indicates knowledge matters more than average in environment $j$
    \item \textbf{No confounding}: Prevents confounding between slope adjustments and the environment main effect $\alpha_j$
\end{enumerate}

\subsubsection{Interpretation of Effective Slopes}
The effective slope $\lambda_{ij}^\text{eff} = \lambda_\text{base} + \theta_j$ determines how strongly model $i$'s knowledge ability $\theta_{ij}^K$ influences success in environment $j$:
\begin{itemize}
    \item If $\theta_j > 0$: Knowledge-intensive environment (knowledge matters more than average)
    \item If $\theta_j < 0$: Execution-heavy environment (knowledge matters less; execution dominates)
    \item If $\theta_j \approx 0$: Knowledge effect is close to the global average
\end{itemize}

Similarly for reasoning: $\psi_{ij}^\text{eff} = \psi_\text{base} + \phi_j$.

\subsubsection{Hierarchical Task Effects}
The task random effects $\eta_t \sim \mathcal{N}(0, \sigma_\eta)$ capture task-specific characteristics not explained by other predictors. The hierarchical prior allows:
\begin{enumerate}
    \item \textbf{Shrinkage}: Tasks with few observations are pulled toward zero (borrowing strength across tasks)
    \item \textbf{Flexibility}: Tasks with strong unique characteristics can deviate substantially from the mean
    \item \textbf{Scale estimation}: $\sigma_\eta$ quantifies how much task-specific variation exists beyond fixed effects
\end{enumerate}

\subsection{Model Performance}

Model 3 achieves excellent predictive performance on held-out data:
\begin{itemize}
    \item \textbf{Task-level LOO correlation}: $r = 0.913$ (explains 83.4\% of variance in task success rates)
    \item \textbf{Environment-level calibration}: $r = 1.000$ (perfect environment-level predictions)
    \item \textbf{Convergence}: All $\hat{R} < 1.01$, no divergences with 4 chains $\times$ 2000 draws
\end{itemize}

\subsection{Variance Decomposition}

The relative importance of each model component (measured as percentage of explained variance on the logit scale):
\begin{itemize}
    \item Reasoning ability: 46.8\%
    \item Environment effects: 34.7\%
    \item Task random effects: 12.8\%
    \item Difficulty level: 3.0\%
    \item Knowledge ability: 1.6\%
    \item Scaffold type: 1.1\%
    \item Verbosity: 0.1\%
\end{itemize}

This decomposition reveals that \textbf{reasoning ability} and \textbf{environment characteristics} are the dominant factors explaining model performance variation, with task-specific effects contributing a moderate amount of additional variance.

\end{document}
